% Agent 12: Pages 388--391 (PDF pages 24--25)
% Covers: Section 3 (end), Section 4 intro, Section 4.1, Section 4.2 (beginning)

%%% ====================================================================
%%% END OF SECTION 3: THE ORIGIN OF STELLAR BLACK HOLES (continued)
%%% ====================================================================

In particular, both observation and theory suggest that mass loss is very
significant during the late stages of stellar evolution and during the death throes
of massive stars. Mass loss might be so important that black holes are rare
beasts, indeed!

Let us focus attention on mass loss during the death throes, as predicted by
the best numerical calculations to date. But in doing so, let us keep in mind the
very primitive state of the modern computations---for example, their typical
neglect of the effects of rotation.

Modern computations conclude that stars whose masses exceed the
Chandrasekhar limit, $M > 1 \, M_\text{core}$, at the endpoint of their evolution should
explode as supernovae. A supernova explosion can leave behind two remnants: an
expanding gas cloud (the outer part of the original star), and a remnant ``star''.
The remnant ``star'' will be a neutron star if its mass, $M_\text{remnant}$, is less than $\sim 2M_\odot$;\footnote{Here one should not ignore the large mass defect; for massive neutron stars, see Zel'dovich and Novikov (1971).}
but will be a black hole if its mass is greater than $\sim 2 M_\text{core}$.

We summarize in Table~3.1 the results of various calculations of the supernova
process.

In Table~3.1 $M_\text{core}$ is formally the mass of the entire star used in the numerical
calculations. However, all the calculations assumed a homogeneous initial model
(typically polytropic or isothermal), whereas the theory of stellar evolution
predicts a highly inhomogeneous structure for presupernova stars. In particular,
stellar evolution produces a central core which is more or less homogeneous,
surrounded by a huge diffuse envelope of iron which is more or less homogeneous,
surrounded by a huge diffuse envelope of iron, hydrogen, and other
elements. Thus, we must imagine the initial stars of the supernova calculations to
be the central, homogeneous core; and we must regard those calculations as
neglecting the envelope. Such neglect is reasonable when one studies supernova
dynamics, because the core is collapsing and it has not enough time to do
anything while the core is collapsing and reexploding. The envelope will generally
be thrown off, with little expenditure of work, by the reexploding parts of the
core. The mass of the envelope might be twice that of the core---or perhaps only
the same as the core, or perhaps even less. The theory of stellar evolution is far
from definitive on this point. Hence the question marks in column~2 of Table~3.1.

It is clear, from the diverse results shown in Table~3.1, that the theory of
supernova explosions is far from perfect in 1972. Nevertheless, one can form the
following very tentative conclusions. (i)~For stars which approach the ends of their
quasistatic evolution with masses less than $\sim$12 to 30$M_\odot$, the supernova explosion
may produce a neutron star. (ii)~For stars with masses greater than $\sim$12 to 30$M_\odot$,
the explosion may produce a black hole. If this tentative conclusion is correct, then
no more than $\sim$1 per cent of the mass of the Galaxy should be in the form of
black holes today; and new black holes should be created at a rate no greater
than $\sim$0.01 per year.

\begin{table*}[htbp]
\caption{Summary of Numerical Computations of Supernova Explosions}
\label{tab:3.1}
\centering
\begin{tabular}{l c c c c}
\hline\hline
& \multicolumn{2}{c}{Initial conditions} & \multicolumn{2}{c}{Results} \\
\cline{2-3}\cline{4-5}
& $M_\text{core}$ & $M_\text{remnant}$ & & \\
Authors & Principal Factor Regulating & $/ M_\odot$ & $/ M_\odot$ & \\
& the Explosion & & & \\
\hline
& & \multicolumn{1}{c}{$10^{53}$~ergs} & & \\[3pt]
Colgate and & Carbon detonation & & & \\
White (1966) & during collapse & 1.5 & 2 & 1.4 \\[3pt]
Arnett (1966) & Carbon detonation in & 0.1 & 0.87 & 3.5 \\
& quasistatic evolution & 0.1 & 0.86 & 9 \\[3pt]
Arnett (1969) & Carbon detonation at the end of & 0.07 & 4 & 1.4 \\
& degenerate evolution & & & \\[3pt]
Fraley (1968) & Oxygen detonation in & 0.025 & 9.75 & 30 \\
& the envelope during & & & \\
& collapse & & & \\[3pt]
Nadezhin (1969) & Oxygen detonation & 0 & $\lesssim$96 & 33 \\
and Ivanova & $\bigg\{$ & 0 & $\lesssim$4 & 8 \\
(1969) & \begin{tabular}{@{}l@{}} $e$-neutrinos: there is \\ no mass loss \end{tabular} & & & \\[3pt]
Arnett (1969) & The core is opaque & 0.66 & 2.46 & 12 \\
& to emission of & 0.50 & 9.56 & 4 \\
& neutrinos & & & \\[3pt]
Wheeler (1969) & Absorption of neutrinos in the & 0.10 & 4 & \\
and Hansen & envelope is opaque & 2.3 & 0.95 & 2 \\[3pt]
(1969) & to neutrinos & 3 & 0 & $\gtrsim 100$ \\[3pt]
Fraley (1968) & Stellar evolution & 1.9 & 1.8 & 30 \\
& & & 0 & 10 \\
\hline\hline
\end{tabular}

\smallskip
\noindent
$\dagger$~In these calculations the initial core and final remnant have the same masses: $M_0 = M_\text{core}$, and since the models give no mass loss, we must
take $2.8 = M/4M_\odot$ and $3 \times 32 = 96 \, M_\odot$ as the masses of the ``true'' remnants.

\noindent
$\ddagger$~In recent calculations by Bruenn (1972) the thermal explosion of a star of $M/M_\odot$ produced a remnant. (This does not change the conclusions given in the text about the formation of black holes.)
\end{table*}


%%% ====================================================================
%%% SECTION 4: BLACK HOLES IN THE INTERSTELLAR MEDIUM
%%% ====================================================================

\section{Black Holes in the Interstellar Medium}
\label{sec:4}

\subsection{Accretion of Noninteracting Particles onto a Nonmoving Black Hole}
\label{sec:4.1}

Consider a black hole at rest in the interstellar medium; and temporarily treat
the interstellar gas as though it were made up of noninteracting particles
[collisions neglected; mean free paths large compared to the region over which
the hole's gravity makes itself felt; $l_\text{fp} \gg 2GM/c^2$; see below]. In the case of a
Schwarzschild hole, all particles with angular momentum per unit mass
$\ell < 2r_g c$  eventually get captured by the hole (see Box~25.6 of MTW or Figure~12
of ZN). Here $r_g = 2GM/c^2$ is the gravitational radius of hole and $M$ is the hole's
mass. If the particle speeds far from the hole are $v_\infty$, then this condition for
capture corresponds to an impact parameter
%
\begin{equation}
b = \ell / v_\infty = 2r_g c/v_\infty\,.
\tag{4.1.1}
\end{equation}
%
Consequently, the ``capture cross section'' of the hole is
%
\begin{equation}
\sigma_\text{capture} = \pi b^2 = 4\pi r_g^2 (c/v_\infty)^2\,.
\tag{4.1.2}
\end{equation}
%
For a Kerr hole the capture cross section is of this same order of magnitude. The
rest mass per unit time crossing inward through a sphere of $r \gg r_\text{capture}$, with
particles directed into the capture region (solid angle $\Delta\Omega = \sigma/r^2$), is
%
\begin{equation}
\dot{M}_0 = \left(\frac{\rho_\text{cap.}}{4\pi}\right) 4\pi r^2 \Delta\Omega = 4\pi r_g^2 \rho_\infty c^2 / v_\infty\,.
\tag{4.1.3}
\end{equation}
%
This is the rate at which the hole accretes rest mass. Rewritten in typical astro\-nomical units, this accretion rate is
%
\begin{equation}
\dot{M}_0 = 10^{11} \left(\frac{\rho_\infty}{10^{-24}~\text{g~cm}^{-3}}\right) \left(\frac{M}{M_\odot}\right)^2 \left(\frac{v_\infty}{10~\text{km~sec}^{-1}}\right)^{-1} \text{``H~II''}
\tag{4.1.3$'$}
\end{equation}
%
\begin{equation}
\frac{d(M_0/M_\odot)}{d(t/10^{10}~\text{yrs})} = 10^{-13} \left(\frac{\rho_\infty}{10^{-24}~\text{g~cm}^{-3}}\right) \left(\frac{M}{M_\odot}\right)^2 \left(\frac{v_\infty}{10~\text{km~sec}^{-1}}\right)^{-1}\,.
\notag
\end{equation}

[The typical densities and speeds of interstellar gas particles in ionized, ``H~II''
regions are $10^{-24}$~g~cm$^{-3}$ and 10~km~sec$^{-1}$; see e.g.\ Kaplan and Pikel'ner (1970)].
Even if 100 per cent of the inflowing rest mass were somehow converted into
outgoing radiation, the total luminosity would be only
%
\begin{equation}
L_\text{max} = \left(10^{24}~\frac{\text{erg}}{\text{sec}}\right) \left(\frac{\rho_\infty}{10^{-24}~\text{g~cm}^{-3}}\right) \left(\frac{M}{M_\odot}\right)^2 \left(\frac{v_\infty}{10~\text{km~sec}^{-1}}\right)^{-1}
\tag{4.1.4}
\end{equation}
%
---a value much too small to be astronomically interesting. Therefore it is fortunate
that the electrons, ions, and magnetic fields of the interstellar gas interact strongly
enough to make the accretion process obey fluid-dynamic laws rather than the
laws of noninteracting particles (see \S4.6, below).

%%% ====================================================================
%%% SECTION 4.2: ADIABATIC, HYDRODYNAMIC ACCRETION
%%% ====================================================================

\subsection{Adiabatic, Hydrodynamic Accretion onto a Nonmoving Black Hole}
\label{sec:4.2}

Switch, then, from a noninteracting description of accretion to a hydrodynamic
description. Assume, as above, that the black hole is at rest with respect to the gas
and a Schwarzschild hole, so that the accretion rate $\dot{M}_0$ and all other basic characteristics of the
hydrodynamic case the flow is spherically symmetric. In the
flow are governed by the gravitational field at distances much greater than the
gravitational radius. (This is because the flow at small radii is supersonic and
therefore cannot influence conditions at large radii.) Thus, one can calculate
the mass flow and other quantities at large radii accurately using the Newtonian
theory of gravitation. Phenomena close to the gravitational radius will be treated
later. The motion of the gas will be assumed adiabatic. Deviations from adiabatic
flow due to radiative losses and radiative transport between gas elements can be
taken into account later by suitably modifying the adiabatic index $\Gamma$.
The form of the flow is governed by two fundamental equations: conservation
of rest mass (``continuity equation''), which we write in the form\footnote{Material for this section is drawn largely from Chapter~13 of ZN, and from Schwarzman and Novikov (1971).}
%
\begin{equation}
4\pi r^2 \rho u = \dot{M}_0 = \text{constant, independent of}~r,
\tag{4.2.1}
\end{equation}
%
and the Euler equation
%
\begin{equation}
u\frac{du}{dr} = -\frac{1}{\rho}\frac{dp}{dr} - \frac{GM}{r^2}\,.
\tag{4.2.2}
\end{equation}
%
[Cf.\ eqs.\ (2.5.23) and (2.5.59).]
Here $\dot{M}_0$ is the total rate of accretion of rest mass, $u$ is the radial velocity, and
$M$ is the mass of the hole. We assume the gas has constant adiabatic index
$\Gamma$, so that during the accretion the pressure and density are related by the
adiabatic law $p = K\rho^\Gamma$, and the speed of sound is given by $a = (\Gamma p/\rho)^{1/2}$. Let
the constants $K$ and $\Gamma$ be given. Our task is to determine the accretion rate $\dot{M}_0$
and to compute as well the distributions of density $\rho(r)$ and velocity $u(r)$ in the
flow.\footnote{Notice that our notation differs from that in \S2. Throughout \S4 we denote (for ease of
eyesight)

(speed of sound) $= a$, not $c_s$

(density of rest mass) $= \rho$, not $\rho_0$

No confusion is likely, since nowhere in \S4 shall we deal with 4-accelerations or with
total mass-energy; and $\rho = \rho_0(1 + \pi)$, and because $\Pi \ll 1$ in all regions of the accretion gas.
We shall also retain factors of $G$, $c$, and $k$ (i.e., use cgs units) throughout \S4.}

In place of the Euler equation (4.2.2) we shall use the Bernoulli equation
(2.5.38) rewritten in the form
%
\begin{equation}
\frac{1}{2} u^2 + \frac{1}{\Gamma - 1} a^2 - \frac{GM}{r} = \text{constant} = \frac{1}{2} u_\infty^2 + \frac{1}{\Gamma - 1} a_\infty^2\,.
\tag{4.2.3}
\end{equation}
%
[Here the enthalpy has been written in the form $w = \int \rho^{-1}\,dp = a^2/(\Gamma - 1)$;
cf.\ eq.\ (2.5.42).] The constant has been determined by conditions at infinity,
where the gas is at rest.

Rewrite the law of mass conservation (4.2.1) with density expressed in terms
of sound speed
%
\begin{equation}
u = \frac{\dot{M}_0}{4\pi \rho_\infty r^2} \left(\frac{a_\infty}{a}\right)^{2/(\Gamma-1)}\,.
\tag{4.2.4}
\end{equation}

The only unknown parameter in the coupled equations (4.2.3) and (4.2.4) is
the accretion rate $\dot{M}_0$. Once $\dot{M}_0$ is known, one can readily solve for the radial
distributions of velocity, sound speed, and density.

The mass flux is determined in the following way. Consider the system of
equations (4.2.3) and (4.2.4). In the $u$--$a$ plane, the Bernoulli equation (4.2.3)
for fixed radius $r$ defines an ellipse; each value of $r$ corresponds to a different ellipse.
See Figure~4.2.1. Similarly, the equation of mass conservation (4.2.4) for each
value of $r$ defines a hyperbola of fractional power, $ua^{2/(\Gamma-1)} = \text{const.}$\ Now, let
the accretion rate $\dot{M}_0$ be chosen arbitrarily. For every value of $r$ (4.2.3) and
(4.2.4) are 2 equations with 2 unknowns, $u$ and $a$. Solving these equations---i.e.,
finding the intersection point of the ellipse with the hyperbola---gives $u$ and $a$
for that particular radius. In other words, in the $u$, $a$ plane the curve $u(a)$ is
determined parametrically by the intersections of corresponding ellipses and
hyperbolae. It is clear (see Figure~4.2.1) that for every ellipse-hyperbola pair,
there are either two intersection points, or one point of tangency, or no
intersection at all. If, for a particular chosen $\dot{M}_0$, there exists any $r$ at which the
curves do not intersect, then the two equations are incompatible for the $\dot{M}_0$, and
such flow cannot occur. Rejecting this case, we have 2 remaining possibilities, for
a given $\dot{M}_0$ (see Figure~4.2.2):

1. The ellipse and hyperbola intersect twice for every value of $r$ (Figure~4.2.2a),
except one $r$, corresponding to 2 families of intersection points.

2. The ellipse and hyperbola meet tangentially (Figure~4.2.2b). In this case the two
families of intersection points $u(a)$ cross each other at $r = r_s$. We shall call the two
(call it $r_s$), at which they meet tangentially the ``bisectrix'' of the graph,
below that this crossing point necessarily lies on the
$u = a$. In both the cases, 1 and 2, the curves $u(a)$ describe possible flows of gas
in the gravitational field. The curves begin on the innermost ellipse, $r = \infty$. On
the upper branch $u(a)$, at this ellipse we have $a = 0$, but $u \neq 0$. These boundary
conditions are not of interest for the accretion problem\footnote{This curve describes the ``stellar wind'' by which some stars eject matter into interstellar
space; see Chapter~13 of ZN.} because accretion
requires $u_\infty = 0$, $a_\infty \neq 0$, i.e., the lower curve $u(a)$ begins at the point $u_\infty = 0$, $a_\infty \neq 0$

\begin{figure}[htbp]
\centering
% \includegraphics[width=\columnwidth]{fig_4_2_1}
\fbox{\parbox{0.9\columnwidth}{\centering [Figure 4.2.1 placeholder]\\
The $u$, $a$ ``solution plane'' for solving the coupled Bernoulli equation (4.2.3)
and law of mass conservation (4.2.4).}}
\caption{The $u$, $a$ ``solution plane'' for solving the coupled Bernoulli equation (4.2.3) and law of mass conservation (4.2.4).}
\label{fig:4.2.1}
\end{figure}

\begin{figure}[htbp]
\centering
% \includegraphics[width=\columnwidth]{fig_4_2_2}
\fbox{\parbox{0.9\columnwidth}{\centering [Figure 4.2.2 placeholder]\\
Possible forms of adiabatic gas flow in a spherical, Newtonian gravitational
field. Curves~(a) [case~1 in text] correspond to flow that remains subsonic or supersonic
everywhere. Curves~(b) [case~2 in text] correspond to flow for which there is a sonic point,
$u = a$.}}
\caption{Possible forms of adiabatic gas flow in a spherical, Newtonian gravitational
field. Curves~(a) [case~1 in text] correspond to flow that remains subsonic or supersonic
everywhere. Curves~(b) [case~2 in text] correspond to flow for which there is a sonic point,
$u = a$.}
\label{fig:4.2.2}
\end{figure}
